\section{Discussion}\label{sec:discussion}

\subsection{Mechanisms}

The evidence describes a minimum wage with little purchase on the pay of the workers it is
meant to protect. The wage null is my most robust result: it holds across two reforms, across
the wage distribution, and across demographic groups. Two features of the institutional
environment make it intelligible. The first is inflation. The 2022 increase of ten percent in
nominal terms met consumer price inflation of roughly eight percent, so the real value of the
floor barely rose and then eroded. The second, and more fundamental, is informality. Most
low-skilled Peruvian workers are informal or self-employed, and for them the statutory minimum
is not enforced; compliance is high only within registered formal firms, which employ a
minority of the low-skilled. A floor that reaches only a small and selected part of the
relevant labor force cannot move its average wage by much, and the earnings distribution in
Figure~\ref{fig:dist} confirms that it did not.

If the reform did not raise pay, did it move anything else? Around 2022 I find a shift from
formal employment toward self-employment, concentrated among the young and among women. The
natural mechanism is the classic one, a binding floor in the covered sector inducing some
reallocation to the uncovered sector, the displacement channel that \citet{jaramillo_lopez_2006}
documents for Peru and \citet{ham_2018_honduras} for Honduras; the accompanying rise in hours
is consistent with movement toward self-employment, where hours are less constrained. I read
this cautiously, because the compositional effect does not survive my most demanding inference
and does not reappear around the 2025 reform. Its reversal out of sample suggests that part of
what I capture around 2022 reflects the uneven post-pandemic recovery rather than a structural
response to the minimum wage, so I present the reallocation as a feature of the 2022 episode
rather than as a general law.

\subsection{Reconciliation with the previous literature}

My results sit comfortably with some of the prior evidence and in tension with the rest, and
the pattern is itself informative. They echo the older Peruvian work of
\citet{jaramillo_lopez_2006}, who found the 2003 reform largely innocuous with displacement
toward informality, and of \citet{cespedes_2005}, who found effects concentrated among the
young. They are consistent with the broad international finding that minimum wage increases in
this range produce small employment effects \citep{card_krueger_1995, cengiz_2019}.

The tension is with the most recent Peruvian study. \citet{jimenez_2023} finds that the 2016
reform raised formal wages by four to five percent and increased formality, a lighthouse
pattern that I do not reproduce for 2022. Three differences plausibly account for the
divergence. The 2016 increase occurred in a low-inflation period, so its real bite was larger
and more durable than that of the inflation-eroded 2022 increase. The identification differs:
\citet{jimenez_2023} uses a department-level dose defined on the formal wage distribution,
while I contrast low- and high-education workers at the individual level, a design that places
more weight on the large informal segment of the low-skilled. And the macroeconomic setting
differs, stable growth in 2016 against a post-pandemic recovery in 2022. Taken together, these
suggest that the labor-market effects of the Peruvian minimum wage are state-dependent, closer
to the lighthouse pattern when the real increase is substantial and the economy is stable, and
muted when inflation erodes the increase and informality absorbs the adjustment.

\subsection{External validity and limitations}

Four limitations bound the interpretation. First, identification rests on the assumption that,
absent the reform, low- and high-education workers would have followed parallel paths. I test
this where possible and bound the consequences of its failure, but I cannot rule out
differential shocks correlated with education in the volatile 2021--2022 window; this is why I
lean on the wage result, whose pre-trends are flat and which replicates out of sample, and
treat the employment-level result as inconclusive. Second, the design identifies effects on
low-skilled workers relative to high-skilled workers. If the reform also reached the
high-skilled, through general-equilibrium spillovers or ripple effects up the distribution, my
estimates mismeasure the absolute effect on the low-skilled; I judge large spillovers to the
tertiary-educated unlikely given that their median wage sits well above the floor, and I note
that because the comparison group is not entirely unexposed, any resulting bias attenuates my
estimates toward zero, so the wage null is a conservative reading. Third, the data are repeated
cross-sections, so I cannot follow individuals across the reform or separate changes in the
composition of the employed from changes in individual outcomes. Fourth, informality in 2024
is measured with a reconstructed indicator rather than the official one; my robustness checks
suggest this does not drive the results, but it adds measurement error. Finally, the window
spans at most eleven post-reform quarters, so I speak to short- and medium-run effects and not
to the long-run adjustment of capital or firm entry and exit.

\subsection{Policy implications}

The results carry a sober message for minimum wage policy in a high-informality economy. A
floor that goes unenforced for most of its intended beneficiaries does little to raise their
pay, and where it does bind, in the formal sector, it may nudge some workers toward less
protected arrangements. Raising the minimum wage is therefore an imperfect instrument for
lifting the earnings of the low-skilled in Peru, a conclusion that echoes
\citet{jaramillo_lopez_2006}. This is not to say the minimum wage is harmless or useless; my
estimates are relative and my window is short. But it does imply that formalization and
enforcement are preconditions for the floor to reach the workers it is designed to protect.
Absent them, the margin of adjustment is the composition of employment rather than the level of
pay.
